
\documentclass[12pt,a4paper]{article}
\usepackage[T2A]{fontenc}
\usepackage[utf8]{inputenc}
\usepackage[russian]{babel}
\usepackage{geometry}
\geometry{top=2.5cm, bottom=2.5cm, left=2cm, right=2cm}
\usepackage{amsmath, amssymb, amsfonts}
\usepackage{tcolorbox}
\tcbuselibrary{breakable}
\usepackage{xcolor}
\usepackage{hyperref}

\definecolor{examplebg}{HTML}{F5F5F5}
\definecolor{exampleframe}{HTML}{CCCCCC}

\newtcolorbox{examplebox}[1][]{%
  breakable, 
  colback=examplebg, 
  colframe=exampleframe, 
  title=Пример, 
  fonttitle=\bfseries, 
  #1
}

\begin{document}

\large

\section{Вариант 1}

\subsection{Задача 1}

\textbf{Описание:}

Вы исследуете заброшенную шахту на хардкорном сервере Майнкрафта. Внезапно ваши факелы закончились, и вы оказались в кромешной темноте. Вы нащупали старый сундук вагонетки, в котором лежат ценные предметы: 15 кучек красной пыли (редстоуна), 12 кристаллов лазурита и 9 изумрудов. На ощупь в темноте все эти предметы абсолютно одинаковые. 

Вам нужно срочно вернуться на базу, чтобы поторговаться с жителем, и для этого нужны именно изумруды. Какое минимальное количество предметов вам нужно вытащить из сундука вслепую, чтобы среди них гарантированно оказалось хотя бы 3 изумруда?

\newpage

\subsection{Задача 2}

\textbf{Описание:}

В Нижнем мире (Незере) на сервере построено 13 порталов, ведущих на разные базы игроков. Главный администратор решил перестроить транспортную сеть и поручил строителям связать порталы скоростными дорогами из синего льда по новому правилу: каждый из 13 порталов должен быть соединен прямой ледовой дорогой ровно с 5 другими порталами. Ледовые дороги могут пересекаться друг с другом на разных высотах, не прерываясь.

Спустя месяц строители заявили, что выполнить это правило в Майнкрафте невозможно, как ни старайся. Правы ли строители? Обоснуйте свой ответ.

\textit{В бланк ответов запишите "Да" или "Нет", а затем подробно распишите ваше доказательство.}

\newpage

\subsection*{Задача 3}

\textbf{Описание:}

В древнем городе под землей находится ровная площадка из $100$ блоков обычного камня, образующая квадрат $10 \times 10$. 
Каждый блок камня соприкасается со своими соседями (сверху, снизу, слева и справа; по диагонали касания нет).

В город пробрался игрок и успел заразить скалком (Sculk) ровно $9$ любых блоков камня на свой выбор, после чего его прогнал Варден. 
Однако скалк обладает свойством автономного разрастания: каждую минуту, если какой-то блок обычного камня соприкасается как минимум с двумя уже зараженными скалком блоками, он тоже превращается в скалк. 

Заражение распространяется до тех пор, пока это возможно. 
Сможет ли игрок так выбрать начальные $9$ блоков, чтобы в итоге скалк поглотил всю площадку (все $100$ блоков)?

\textit{В бланк ответов запишите «Да» или «Нет», а затем строго обоснуйте свой ответ.}

\newpage
\subsection*{Задача 4}

\textbf{Описание:}

В гигантском механизме из красного камня (редстоуна) установлена длинная линия из ровно $5000$ редстоуновых ламп. Все лампы пронумерованы по порядку от $1$ до $5000$.
Изначально все лампы находятся в выключенном состоянии (режим «0»).

Каждую игровую ночь командный блок запускает алгоритм «Каскадного переключения», который состоит ровно из $5000$ тиков (шагов).
\begin{itemize}
    \item В Тик 1 командный блок меняет состояние у каждой лампы (переключает все лампы из «Выкл» в «Вкл»).
    \item В Тик 2 система меняет состояние у каждой второй лампы (то есть у ламп 2, 4, 6, 8...). Если лампа не горела — она загорается, если горела — гаснет.
    \item \dots
    \item В Тик $K$ система переключает состояние у каждой $K$-й лампы.
    \item Наконец, в Тик 5000 система меняет состояние только у одной, $5000$-й лампы.
\end{itemize}

После завершения всех $5000$ тиков механизм останавливается. Все лампы, которые остались во включенном состоянии (режим «1»), используются для освещения главной базы.

\textbf{Требуется ответить на два вопроса:}
\begin{enumerate}
    \item Какое общее количество ламп останется во включенном состоянии в итоге?
    \item Какой номер будет у самой последней (наибольшей по номеру) включенной лампы?
\end{enumerate}

\textit{В бланк ответов запишите два числа через пробел. Затем приведите обоснование для каждого из ответов.}

\newpage

\subsection*{Задача 5}

\textbf{Описание:}

Вы построили гигантскую плоскую арену, обнесенную стенами из блоков слизи. 
Её внутренние размеры: ровно $1920$ блоков в ширину и $1080$ блоков в длину.

В четырех углах этой арены установлены блоки-мишени (Target blocks), которые улавливают снаряды. Координаты углов арены:
\begin{itemize}
    \large
    \item $(0, 0)$ — Левый нижний угол (точка старта)
    \item $(1920, 0)$ — Правый нижний угол
    \item $(0, 1080)$ — Левый верхний угол
    \item $(1920, 1080)$ — Правый верхний угол
\end{itemize}

Из Раздатчика, установленного в Левом нижнем углу $(0,0)$, выстреливают зачарованной стрелой ровно под углом $45^\circ$ (строго по диагонали). 
Стрела летит по идеальной прямой, игнорируя гравитацию. Каждый раз, когда она врезается в стену из блоков слизи (вертикальную или горизонтальную), она идеально отскакивает от неё (угол падения равен углу отражения) и продолжает лететь по диагонали, но в другую сторону. 

Стрела будет летать по арене до тех пор, пока не попадет точно в один из четырех блоков-мишеней. Как только стрела попадает в мишень, она вонзается в нее, и полет заканчивается.

\textbf{Требуется определить:}
\begin{enumerate}
    \item В какую именно из трех оставшихся мишеней попадет стрела в конце своего пути?
    \item Какое суммарное количество раз стрела отскочит от стен из слизи перед тем, как попасть в мишень? (Само финальное попадание в мишень отскоком от стены не считается).
\end{enumerate}

\textit{В бланк ответов запишите координаты угла (мишени) и одно целое число — количество отскоков. Далее, запишите ваше решение.}

\end{document}
